\documentclass[10pt,a4paper]{amsart}
\usepackage{amsmath,amssymb,amsthm,graphicx}
% Retained preamble for the Stacks Project chapter morphisms.tex (Morphisms of Schemes).
% The chapter file pinned-original-source.tex carries only \input{preamble} on its first line, so no
% preamble text lives inside the chapter itself.  The Stacks Project uses one single shared preamble for
% every chapter; the snapshot stages it verbatim in sources/stacks-chow/preamble.tex.  The block below is
% copied mechanically, line for line, from that file's lines 14-18 (the Xy-pic setup) and lines 154-217
% (the native theorem declarations together with the mathematical macros).  Nothing was added, removed or
% redefined.  The wrapper class and the amsmath/amssymb/amsthm/graphicx base packages are supplied by the
% assembler, so the class/\externaldocument part of the shared preamble is not repeated here.
% For commutative diagrams we use Xy-pic
\usepackage[all]{xy}

% We use 2cell for 2-commutative diagrams.
\xyoption{2cell}
% Theorem environments.
%
\theoremstyle{plain}
\newtheorem{theorem}[subsection]{Theorem}
\newtheorem{proposition}[subsection]{Proposition}
\newtheorem{lemma}[subsection]{Lemma}

\theoremstyle{definition}
\newtheorem{definition}[subsection]{Definition}
\newtheorem{example}[subsection]{Example}
\newtheorem{exercise}[subsection]{Exercise}
\newtheorem{situation}[subsection]{Situation}

\theoremstyle{remark}
\newtheorem{remark}[subsection]{Remark}
\newtheorem{remarks}[subsection]{Remarks}

\numberwithin{equation}{subsection}

% Macros
%
\def\lim{\mathop{\mathrm{lim}}\nolimits}
\def\colim{\mathop{\mathrm{colim}}\nolimits}
\def\Spec{\mathop{\mathrm{Spec}}}
\def\Hom{\mathop{\mathrm{Hom}}\nolimits}
\def\Ext{\mathop{\mathrm{Ext}}\nolimits}
\def\SheafHom{\mathop{\mathcal{H}\!\mathit{om}}\nolimits}
\def\SheafExt{\mathop{\mathcal{E}\!\mathit{xt}}\nolimits}
\def\Sch{\mathit{Sch}}
\def\Mor{\mathop{\mathrm{Mor}}\nolimits}
\def\Ob{\mathop{\mathrm{Ob}}\nolimits}
\def\Sh{\mathop{\mathit{Sh}}\nolimits}
\def\NL{\mathop{N\!L}\nolimits}
\def\CH{\mathop{\mathrm{CH}}\nolimits}
\def\proetale{{pro\text{-}\acute{e}tale}}
\def\etale{{\acute{e}tale}}
\def\QCoh{\mathit{QCoh}}
\def\Ker{\mathop{\mathrm{Ker}}}
\def\Im{\mathop{\mathrm{Im}}}
\def\Coker{\mathop{\mathrm{Coker}}}
\def\Coim{\mathop{\mathrm{Coim}}}

% Boxtimes
%
\DeclareMathSymbol{\boxtimes}{\mathbin}{AMSa}{"02}

%
% Macros for moduli stacks/spaces
%
\def\QCohstack{\mathcal{QC}\!\mathit{oh}}
\def\Cohstack{\mathcal{C}\!\mathit{oh}}
\def\Spacesstack{\mathcal{S}\!\mathit{paces}}
\def\Quotfunctor{\mathrm{Quot}}
\def\Hilbfunctor{\mathrm{Hilb}}
\def\Curvesstack{\mathcal{C}\!\mathit{urves}}
\def\Polarizedstack{\mathcal{P}\!\mathit{olarized}}
\def\Complexesstack{\mathcal{C}\!\mathit{omplexes}}
% \Pic is the operator that assigns to X its picard group, usage \Pic(X)
% \Picardstack_{X/B} denotes the Picard stack of X over B
% \Picardfunctor_{X/B} denotes the Picard functor of X over B
\def\Pic{\mathop{\mathrm{Pic}}\nolimits}
\def\Picardstack{\mathcal{P}\!\mathit{ic}}
\def\Picardfunctor{\mathrm{Pic}}
\def\Deformationcategory{\mathcal{D}\!\mathit{ef}}

% --- appended 2026-10-12 (Chief): the project own definitions of the reference/slogan/history
% display environments, copied mechanically from the shared Stacks preamble lines 8-12
% (\usepackage{verbatim} plus the three \newenvironment lines, which the chapter preamble keeps
% verbatim from the project preamble). Nothing added, removed or redefined: these environments
% carry no mathematical content and the project itself typesets them as comments. Without these
% three lines a counted statement whose own text uses them fails with Environment slogan undefined.
\usepackage{verbatim}
\newenvironment{reference}{\comment}{\endcomment}
\newenvironment{slogan}{\comment}{\endcomment}
\newenvironment{history}{\comment}{\endcomment}
\usepackage{cleveref}
\crefname{originallemma}{original lemma}{original lemmas}
\Crefname{originallemma}{Original lemma}{Original lemmas}
\makeatletter
\newlabel{morphisms-lemma-set-points-where-fibres-etale}{{0476}{source}{}{sourcepointer-1}{}}
\newlabel{morphisms-lemma-set-points-where-fibres-etale@cref}{{[originallemma][0][]0476}{[1][0][]source}{}{}{}}
\makeatother
\crefname{originallemma}{original lemma}{original lemmas}
\Crefname{originallemma}{Original lemma}{Original lemmas}
\makeatletter
\newlabel{lemma-set-points-where-fibres-unramified}{{0475}{source}{}{sourcepointer-2}{}}
\newlabel{lemma-set-points-where-fibres-unramified@cref}{{[originallemma][0][]0475}{[1][0][]source}{}{}{}}
\makeatother
\crefname{originallemma}{original lemma}{original lemmas}
\Crefname{originallemma}{Original lemma}{Original lemmas}
\makeatletter
\newlabel{lemma-flat-unramified-etale}{{02GV}{source}{}{sourcepointer-3}{}}
\newlabel{lemma-flat-unramified-etale@cref}{{[originallemma][0][]02GV}{[1][0][]source}{}{}{}}
\makeatother
\crefname{originallemma}{original lemma}{original lemmas}
\Crefname{originallemma}{Original lemma}{Original lemmas}
\makeatletter
\newlabel{lemma-base-change-module-flat}{{01U8}{source}{}{sourcepointer-4}{}}
\newlabel{lemma-base-change-module-flat@cref}{{[originallemma][0][]01U8}{[1][0][]source}{}{}{}}
\makeatother
\crefname{originallemma}{original lemma}{original lemmas}
\Crefname{originallemma}{Original lemma}{Original lemmas}
\makeatletter
\newlabel{lemma-composition-module-flat}{{01U6}{source}{}{sourcepointer-5}{}}
\newlabel{lemma-composition-module-flat@cref}{{[originallemma][0][]01U6}{[1][0][]source}{}{}{}}
\makeatother
\crefname{originallemma}{original lemma}{original lemmas}
\Crefname{originallemma}{Original lemma}{Original lemmas}
\makeatletter
\newlabel{lemma-flat-permanence}{{02JZ}{source}{}{sourcepointer-6}{}}
\newlabel{lemma-flat-permanence@cref}{{[originallemma][0][]02JZ}{[1][0][]source}{}{}{}}
\makeatother
\crefname{originallemma}{original lemma}{original lemmas}
\Crefname{originallemma}{Original lemma}{Original lemmas}
\makeatletter
\newlabel{lemma-base-change-smooth}{{01VB}{source}{}{sourcepointer-7}{}}
\newlabel{lemma-base-change-smooth@cref}{{[originallemma][0][]01VB}{[1][0][]source}{}{}{}}
\makeatother
\crefname{originallemma}{original lemma}{original lemmas}
\Crefname{originallemma}{Original lemma}{Original lemmas}
\makeatletter
\newlabel{lemma-etale-at-point}{{02GU}{source}{}{sourcepointer-8}{}}
\newlabel{lemma-etale-at-point@cref}{{[originallemma][0][]02GU}{[1][0][]source}{}{}{}}
\makeatother
\title{The etale locus is compatible with base change under flatness or finite presentation hypotheses}
\author{The Stacks Project Authors}
\date{stacks-project:morphisms.tex:local-snapshot-20261012; GNU FDL 1.2 or later}
\begin{document}
\maketitle
\setcounter{subsection}{1}
\noindent\textit{Editorial note.} All delivered mathematical text is the author's own, copied from the pinned Stacks Project chapter file morphisms.tex; no statement or proof step was rewritten, reordered or removed.\par\medskip
\noindent\textit{Editorial note.} The chapter file pinned-original-source.tex carries only \textbackslash{}input\{preamble\} on its first line, so no preamble text lives inside the chapter. The retained file sources/stacks-morphisms/retained-preamble.tex is a mechanical, line-for-line copy of the project's single shared preamble as staged in the snapshot. Nothing was added, removed or redefined, and the excerpt serves only as preamble, never as a span coordinate.\par\medskip
\noindent\textit{Editorial note.} Both counted spans start at the \textbackslash{}begin line of their environment, so the statement and the author's proof are delivered whole; surrounding prose of the section is deliberately outside the counted spans and is carried as standard prerequisites instead.\par\medskip
\noindent\textit{Editorial note.} The citation fields (source\_id, source\_version, authors, source\_url, license.name, source\_title) are identical to those of every other card this lane has produced from this chapter, so the batch selector treats them as one source.\par\medskip
\noindent\textit{Allowed standard prerequisites (editorial).}\begin{itemize}
\item The unramified locus is compatible with base change. Lemma \textbackslash{}ref\{lemma-set-points-where-fibres-unramified\}, cited at proof lines 8434-8440. Gives the equality of the two unramified point sets.
\item A morphism is etale at a point exactly when it is flat and unramified there. Lemma \textbackslash{}ref\{lemma-flat-unramified-etale\}, cited at proof lines 8443-8444. Turns the unramified statement into the etale statement in case (1).
\item Flatness is preserved by base change and by composition, and is detected at the source point. Lemma \textbackslash{}ref\{lemma-base-change-module-flat\}, Lemma \textbackslash{}ref\{lemma-composition-module-flat\} and Lemma \textbackslash{}ref\{lemma-flat-permanence\}, cited at proof lines 8447-8453. Transports flatness in case (2).
\item The etale-at-a-point criterion. Lemma \textbackslash{}ref\{lemma-etale-at-point\}, cited at proof lines 8454-8456. Concludes that the morphism is etale at the image point in case (2).
\end{itemize}\medskip
\ifdefined\hypertarget\hypertarget{sourcepointer-1}{}\fi
\noindent\textit{Source reference.} Original lemma 0476: focused-seeds/AGP/sources/stacks-snapshot-20261011/tags.txt, source lines 5497--5497. The lemma's own canonical tag, quoted so that readers can resolve the label inside the delivered text; it points at the counted statement itself and carries no proof content.\par\medskip
\ifdefined\hypertarget\hypertarget{sourcepointer-2}{}\fi
\noindent\textit{Source reference.} Original lemma 0475: focused-seeds/AGP/sources/stacks-snapshot-20261011/tags.txt, source lines 5496--5496. Navigation pointer to a statement the counted proof refers to but does not contain; nothing load-bearing is delegated to it.\par\medskip
\ifdefined\hypertarget\hypertarget{sourcepointer-3}{}\fi
\noindent\textit{Source reference.} Original lemma 02GV: focused-seeds/AGP/sources/stacks-snapshot-20261011/tags.txt, source lines 3242--3242. Navigation pointer to a statement the counted proof refers to but does not contain; nothing load-bearing is delegated to it.\par\medskip
\ifdefined\hypertarget\hypertarget{sourcepointer-4}{}\fi
\noindent\textit{Source reference.} Original lemma 01U8: focused-seeds/AGP/sources/stacks-snapshot-20261011/tags.txt, source lines 2401--2401. Navigation pointer to a statement the counted proof refers to but does not contain; nothing load-bearing is delegated to it.\par\medskip
\ifdefined\hypertarget\hypertarget{sourcepointer-5}{}\fi
\noindent\textit{Source reference.} Original lemma 01U6: focused-seeds/AGP/sources/stacks-snapshot-20261011/tags.txt, source lines 2399--2399. Navigation pointer to a statement the counted proof refers to but does not contain; nothing load-bearing is delegated to it.\par\medskip
\ifdefined\hypertarget\hypertarget{sourcepointer-6}{}\fi
\noindent\textit{Source reference.} Original lemma 02JZ: focused-seeds/AGP/sources/stacks-snapshot-20261011/tags.txt, source lines 3354--3354. Navigation pointer to a statement the counted proof refers to but does not contain; nothing load-bearing is delegated to it.\par\medskip
\ifdefined\hypertarget\hypertarget{sourcepointer-7}{}\fi
\noindent\textit{Source reference.} Original lemma 01VB: focused-seeds/AGP/sources/stacks-snapshot-20261011/tags.txt, source lines 2440--2440. Navigation pointer to a statement the counted proof refers to but does not contain; nothing load-bearing is delegated to it.\par\medskip
\ifdefined\hypertarget\hypertarget{sourcepointer-8}{}\fi
\noindent\textit{Source reference.} Original lemma 02GU: focused-seeds/AGP/sources/stacks-snapshot-20261011/tags.txt, source lines 3241--3241. Navigation pointer to a statement the counted proof refers to but does not contain; nothing load-bearing is delegated to it.\par\medskip
\begin{lemma}
\label{lemma-set-points-where-fibres-etale}
Let
$$
\xymatrix{
X' \ar[r]_{g'} \ar[d]_{f'} & X \ar[d]^f \\
S' \ar[r]^g & S
}
$$
be a cartesian diagram of schemes. Let $W \subset X$, resp.\ $W' \subset X'$
be the open subscheme of points where $f$, resp.\ $f'$ is \'etale.
Then $W' = (g')^{-1}(W)$ if
\begin{enumerate}
\item $f$ is flat and locally of finite presentation, or
\item $f$ is locally of finite presentation and $g$ is flat.
\end{enumerate}
\end{lemma}

\begin{proof}
Assume first that $f$ locally of finite type. Consider the set
$$
T = \{x \in X \mid f\text{ is unramified at }x\}
$$
and the corresponding set $T' \subset X'$ for $f'$. Then
$T' = (g')^{-1}(T)$ by
Lemma \ref{lemma-set-points-where-fibres-unramified}.

\medskip\noindent
Thus case (1) follows because in case (1) $T$ is the (open) set of points
where $f$ is \'etale by Lemma \ref{lemma-flat-unramified-etale}.

\medskip\noindent
In case (2) let $x' \in W'$. Then $g'$ is flat at $x'$
(Lemma \ref{lemma-base-change-module-flat}) and
$g \circ f'$ is flat at $x'$ (Lemma \ref{lemma-composition-module-flat}).
It follows that $f$ is flat at $x = g'(x')$
by Lemma \ref{lemma-flat-permanence}. On the other hand, since
$x' \in T'$ (Lemma \ref{lemma-base-change-smooth})
we see that $x \in T$. Hence $f$ is \'etale at $x$ by
Lemma \ref{lemma-etale-at-point}.
\end{proof}

\end{document}
